%!TEX program = lualatex

\documentclass[11pt, a4paper]{article}
\usepackage[a4paper, margin=2.5cm]{geometry}
\usepackage{fontspec}
\usepackage[bulgarian, english, bidi=basic, provide=*]{babel}
\babelprovide[import, onchar=ids fonts]{bulgarian}
\babelfont{rm}{Noto Sans}
\babelfont{tt}{Noto Sans Mono}

\usepackage{xcolor}
\usepackage{hyperref}
\usepackage{booktabs}
\usepackage{enumitem}
\usepackage{listings}
\usepackage{titlesec}
\usepackage{xcolor}
\usepackage{tikz}
\usepackage{array}
\usepackage{caption}
\usepackage{titling}

\setlength{\droptitle}{-2cm}

\titleformat{\title}{\normalfont\Large\bfseries}{\thetitle}{1em}{}
\titleformat{\section}{\normalfont\large\bfseries}{\thesection}{1em}{}

\definecolor{primary}{RGB}{30, 64, 175}
\definecolor{secondary}{RGB}{71, 85, 105}
\definecolor{bgcode}{rgb}{0.93,0.93,0.93}
\definecolor{bordercode}{rgb}{0.82,0.82,0.82}

\hypersetup{
    colorlinks=true,
    linkcolor=primary,
    urlcolor=primary
}

\lstset{
    backgroundcolor=\color{bgcode},
    frame=single,
    rulecolor=\color{bordercode},
    basicstyle=\ttfamily\small,
    breaklines=true,
    showstringspaces=false
}

% Създаване на команда \code{...} за inline код със сив фон:
\usepackage[most]{tcolorbox}
\usepackage{varwidth}

\newtcbox{\code}{%
  on line,
  boxsep=0pt,
  left=2pt,
  right=2pt,
  top=1pt,
  bottom=1pt,
  colback=bgcode,
  colframe=bgcode,
  fontupper=\ttfamily,
  tcbox raise base,
  varwidth upper=\linewidth % Използва varwidth за ограничаване до ширината на реда
}

% Команда за бутон +
\newcommand{\addbutton}{%
  \begin{tikzpicture}[baseline=-0.5ex]
    \node[circle, fill=blue!80!black, inner sep=1.5pt, text=white, font=\bfseries\small] {+};
  \end{tikzpicture}%
}

\setlist[itemize]{label=\textbullet}

\title{\Large\bfseries Инструкции за инсталиране на Jupyter Notebook чрез Anaconda}
\date{}

\begin{document}

\maketitle
\vspace{-3cm}

\section{Инсталация на Anaconda}

{\color{red} Не е нужно ако използвате Google Colab!} \\

\textbf{Anaconda} е пакет, който включва Python, Jupyter Notebook и най-популярните библиотеки за обработка на данни.

\begin{enumerate}
    \item Отидете на официалния уебсайт: \url{https://www.anaconda.com/download}
    \item Регистрирайте се за да получите достъп.
    \item Изтеглете инсталационния файл \textbf{Anaconda distribution} за вашата операционна система (Windows, macOS или Linux).
    \item Стартирайте изтегления файл и следвайте инструкциите на екрана.
    \item След завършване на инсталацията вече разполагате с Python и Jupyter Notebook.
\end{enumerate}

\section{Стартиране на Jupyter Notebook}

Има няколко достъпни варианта за работа с Jupyter Notebook: локален Jupyter Notebook сървър, JupyterLab, VS Code, или Google Colab. За сравнение вижте таблицата по-долу.

\subsection{Jupyter Notebook сървър}

\begin{enumerate}
    \item Отворете Anaconda Navigator.
    \item Кликнете върху иконата \textbf{Launch} под Jupyter Notebook.
    \item Това отваря Jupyter Notebook в папката на вашия потребителски профил.
    \item Отворете съществуваща папка или създайте нова на удобно за вас място.
    \item Създайте нов ноутбук с \addbutton{} бутона.
\end{enumerate}

\subsection{JupyterLab}
\begin{enumerate}
    \item Отворете Anaconda Navigator.
    \item Кликнете върху иконата \textbf{Launch} под JupyterLab.
    \item Отворете съществуваща папка или създайте нова на удобно за вас място.
    \item Създайте нов ноутбук с \addbutton{} бутона.
\end{enumerate}

\subsection{VS Code}
\begin{enumerate}
    \item Отворете Anaconda Navigator.
    \item Кликнете върху иконата \textbf{Install} под VS Code.
    \item След инсталация кликнете върху иконата \textbf{Launch} под VS Code.
    \item Включете автоматичното запаметяване \textbf{File} \rightarrow \ \ \textbf{Auto save}.
    \item Отворете съществуваща папка.
    \item Създайте нов \code{.ipynb} файл.
    \item Отворете командния панел с \texttt{Ctrl + Shift + P} и потърсете \textbf{Python: Select Interpreter}, след което изберете интерпретатора от Anaconda.
\end{enumerate}

\subsection{Google Colab}
\begin{enumerate}
    \item Отидете на уебсайта \url{https://colab.research.google.com}.
    \item Влезте в профила си в Googlе.
    \item В началния прозорец изберете \textbf{New Notebook} или от менюто изберете \textbf{File} $\rightarrow$ \textbf{New notebook}.
    \item При нужда свържете файла с вашия Google Drive чрез менюто \textbf{Files} $\rightarrow$ \textbf{Mount Drive}.
    \item Промените се запазват автоматично в Google Drive.
\end{enumerate}

\section{Инсталиране на библиотеки и управление на средите}

\begin{enumerate}
    \item Потърсете и отворете Anaconda Prompt.
    \item Ще видите че автоматично се активира средата \code{base}.
    \item Тя вече съдържа голяма част от нужните библиотеки. (напр. \code{numpy}, \code{pandas})
    \item За да инсталирате нова библиотека използвайте командата \code{pip install [lib]}. Тук \code{[lib]} се заменя с името на библиотеката.
    \item За да създатете нова среда изполвайте командата \code{conda create --name [env]}. Тук \code{[env]} се заменя с името на новата среда.
    \item По подразбиране новите среди не съдържат всички библиотеки и трябва да се инсталират.
    \item За да инсталирате библиотеки в средата, тя трябва да се активира с командата \code{conda activate [env]}. 
    \item За да ползвате средата е от значение кой от вариантите препдпочитате:
    \begin{itemize}
        \item За Jupyter Notebook първо инсталирайте пакета \code{conda install notebook}. След това стартирайте нов сървър с \code{jupyter notebook}.
        \item За JupyterLab първо инсталирайте пакета \code{conda install jupyterlab}. След това стартирайте нов сървър с \code{jupyter lab}.
        \item За VS Code първо инсталирайте пакета \code{conda install ipykernel}. Отворете VS Code в желаната папка/файл и изберете интерпретатора на новата среда.
    \end{itemize}
\end{enumerate}

\begin{table}[htbp]
\centering
\caption*{Сравнение на варианти за работа с Jupyter Notebook}
\label{tab:jupyter_environments}
\begin{tabular}{p{0.22\textwidth} p{0.36\textwidth} p{0.36\textwidth}}
\toprule
\textbf{Вариант} & \textbf{Предимства} & \textbf{Недостатъци} \\
\midrule
\textbf{Jupyter Notebook сървър} & 
\begin{itemize}[leftmargin=*, nosep, before=\vspace{-0.5\baselineskip}]
    \item Ползва се през браузъра. Няма нужда от инсталиране на програми (освен пакета на Anaconda).
\end{itemize} & 
\begin{itemize}[leftmargin=*, nosep, before=\vspace{-0.5\baselineskip}]
    \item Остарял интерфейс. Няма табове, няма страничен панел с файлово дърво.
\end{itemize} \\
\addlinespace[1em]

\textbf{JupyterLab} & 
\begin{itemize}[leftmargin=*, nosep, before=\vspace{-0.5\baselineskip}]
    \item Може да се използва през браузъра, но и да се инсталира като отделна програма.
    \item Интерфейса наподобява IDE програмите. Позволява табове, разделяне на екрана, има файлова структура, разширения, ...
    \item Има вграден \textit{autocomplete} (автодовършване на функции) при двойно натискане на \code{Tab}.
\end{itemize} & 
\begin{itemize}[leftmargin=*, nosep, before=\vspace{-0.5\baselineskip}]
    \item Не е пълно IDE. По-слаба екосистема от разширения - за Git, Latex, Копилот и пр.
\end{itemize} \\
\addlinespace[1em]

\textbf{VS Code} & 
\begin{itemize}[leftmargin=*, nosep, before=\vspace{-0.5\baselineskip}]
    \item Леко IDE, с богата екосистема и приложения, повечето от които са напълно безплатни.
    \item Покрива всички базисни качества на JupyterLab.
    \item По-добър \textit{autocomplete}. Включва и други предимства от същия тип, като показване на плаващи прозорци с документация на функции, ...
\end{itemize} & 
\begin{itemize}[leftmargin=*, nosep, before=\vspace{-0.5\baselineskip}]
    \item Сложната екосистема понякога води до софтуерни грешки или недобре поддържани разширения.
    \item По-трудно за използване от хора без опит в програмирането.
\end{itemize} \\
\addlinespace[1em]

\textbf{Google Colab} & 
\begin{itemize}[leftmargin=*, nosep, before=\vspace{-0.5\baselineskip}]
    \item Напълно онлайн. Не използвате процесора на локалния си компютър, а на облачна машина.
    \item Възможност за използване на графични процесори за тежки операции.
    \item Не изисква инсталиране на отделни програми на вашия компютър.
    \item Може лесно да се сподели за екипна работа.
\end{itemize} & 
\begin{itemize}[leftmargin=*, nosep, before=\vspace{-0.5\baselineskip}]
    \item Изисква качване на всички данни в интернет, които понякога може да са големи файлове, което отнема време.
\end{itemize} \\
\bottomrule
\end{tabular}
\end{table}


\end{document}